\documentclass[11pt]{article}
\usepackage[margin=1in]{geometry}
\usepackage{amsmath,amssymb}
\usepackage{booktabs}
\usepackage{graphicx}
\usepackage[hidelinks]{hyperref}
\usepackage{float}
\usepackage{setspace}
\usepackage[round]{natbib}
\usepackage[hang,flushmargin]{footmisc}
\onehalfspacing

% threeparttable for tablenotes
\usepackage{threeparttable}

\newcommand{\jk}{Jaroci\'{n}ski--Karadi}

\title{\textbf{What the Fed Reveals and Where Investors Put Their Money}\\[10pt]
\large Fund-Flow Responses to Monetary-Policy and Central-Bank-Information Shocks}
\author{Draft: November 2026\thanks{%
Preliminary and incomplete; comments welcome. All numbers are generated by the
archived pipeline (\texttt{direction2/robustness\_tests.py},
\texttt{professor\_revisions.py}) and are reproducible from the replication
package. We thank the seminar participants at XJTLU for comments.}}
\date{}

\begin{document}
\maketitle

\begin{abstract}
\noindent
Twenty years of evidence establish that Federal Reserve communication moves asset
prices. We show that it also moves \emph{where investors put their money}. Using
the official \jk{} (2020) decomposition of high-frequency announcement surprises
into a conventional monetary-policy (MP) shock and a central-bank-information
(CBI) shock, matched to all 161 FOMC meetings from 2006 to 2026, we relate the
two shock components to monthly mutual-fund flows across seven asset classes.
Three results stand out. First, CBI shocks reallocate portfolios: a one-standard-deviation
favourable-information shock shifts class-month flows up the risk ladder
($+5.28$, $p=.005$ with month-clustered standard errors) and out of government-bond
funds ($-43.1$, $p=.037$), while MP shocks of equal magnitude leave flows
unchanged ($p=.47$--$.95$). Second, the reallocation is persistent but not
permanent: roughly half unwinds two months after the announcement, the
portfolio counterpart of the ``prices first, portfolios later'' pattern in
returns. Third, on the price side, the forward-guidance-period response loads
on the \emph{breadth} of statement sentiment, not on directional hawkishness:
sentiment predicts FG-period Treasury returns ($p=.005$--$.010$) even when both
measures enter together, while the LLM hawkishness score does not ($p=.081$--$.33$).
The flows results survive an identification switch that our own
sign-restriction series fails: the statement-based CBI proxy is nearly
uncorrelated with the official CBI ($\rho=.04$), so the two capture different
information components. (208 words)
\end{abstract}

\newpage

\section{Introduction}\label{sec:intro}

Does the information in a Federal Reserve statement change not only what assets
are worth, but where investors actually put their money? The existing literature
answers the first question emphatically. Since \cite{BERNANKE_KUTTNER_2005}
documented that a few sentences move equity values, a large body of work has
established that the \emph{information} component of monetary-policy news ---
what the statement reveals about the economy, as distinct from the policy action
itself --- is a first-order driver of announcement-day price movements
\citep{JAROCINSKI_KARADI_2020, BAUER_SWANSON_2023}. The second question has
received far less attention, and what evidence exists is indirect
\citep{KROENCKE_2024, BANEGAS_2022, BROOKS_2018}. This paper provides a direct
answer: central-bank-information shocks generate large, statistically robust,
and economically meaningful reallocations of mutual-fund portfolios, while
conventional monetary-policy shocks of the same size leave fund flows
essentially untouched.

\subsection{What we do}\label{sec:whatwedo}

Our identification starts from the \jk{} (2020) insight that announcement-window
co-movement of interest rates and stock prices separates two components of
policy news: an MP shock (rates and stocks move in opposite directions, as a
pure policy action implies) and a CBI shock (they move together, as pure news
about the economy implies). We use Jaroci\'{n}ski's publicly maintained update
of this decomposition, which covers every FOMC announcement through 2026, and
match it to our sample of 161 meetings. On the flow side, we observe
class-month-level fund flows for seven asset classes --- government bonds,
corporate bonds, real assets, large-cap equity, developed-market equity,
emerging-market equity, and small-cap equity --- constructed from CRSP
mutual-fund database holdings and total net assets with a point-in-time
discipline (\S\ref{sec:data}).

The dependent variable in the flow regressions is the class-level flow rate
(\% of lagged TNA) in the meeting month, the following month, and their
difference. The controls are log lagged TNA and 12-month flow volatility; on
the panel specifications, standard errors are clustered by meeting month.

\subsection{What we find}\label{sec:findings}

\paragraph{1. Information shocks reallocate portfolios; policy shocks do not.}
Table~\ref{tab:official} reports the headline. A favourable CBI shock shifts
flows \emph{up} the risk ladder: the Risk$\times$CBI interaction is $+5.28$
($p=.005$, month-clustered) in the meeting month and remains $+3.45$
($p=.006$) one month after. Consistently, government-bond fund flows respond
negatively to the same shock ($-43.1$ per unit of CBI, $p=.037$; the Wald test
rejects equality of MP and CBI coefficients at $p=.025$). The conventional MP
shock --- the same regressor that moves Treasury yields on impact --- produces
nothing anywhere on the risk ladder ($p=.47$--$.95$). Investors do not
realocate in response to the action; they reallocate in response to what the
statement \emph{reveals}.

\paragraph{2. The reallocation unwinds.} In the difference window (post-month
minus same-month flow), the Risk$\times$CBI coefficient falls to $-1.79$
($p=.37$): roughly half of the initial risk-on shift is reversed by month two.
This is the portfolio counterpart of the post-announcement drift and reversal
pattern documented for prices, and it links the flow evidence directly to the
``prices first, portfolios later'' logic of \S\ref{sec:prices}.

\paragraph{3. On the price side, it is breadth, not direction.} The
forward-guidance era (56 meetings, 2008--2015) is where statement language
mattered most for prices. We show that this effect is driven by the
\emph{breadth} of sentiment --- the combined dictionary score that aggregates
many tone dimensions --- rather than by directional hawkishness. In a horse
race with both measures, sentiment survives ($-1.11$, $p=.005$) and the
five-model-consensus LLM hawkishness score does not ($p=.081$--$.33$). The
tone effect also survives the professor-requested horse race with the official
MP and CBI shocks as controls ($-0.96$, $p=.009$; Table~\ref{tab:horse}).

\paragraph{4. Identification discipline.} Our statement-based sign-restriction
series --- the measure used in earlier drafts --- turns out to be a poor proxy
for the official CBI: the two are correlated at $\rho=.04$. Under the
sign-rule series the flow results are fragile (they collapse under HC1 or
clustered standard errors, and hinge on the 2008-09-16 meeting); under the
official series they are robust. We report both throughout, treat the official
series as the primary measure, and view the exercise as a caution for the
growing literature that proxies information shocks from statement text alone.

\subsection{Why it matters}\label{sec:whymatters}

Three implications. First, the flows evidence gives the information effect a
\emph{quantity} dimension: it is not only that prices embed the Fed's private
information, but that investors act on it, moving billions of dollars across
asset classes. Second, the differential response --- CBI yes, MP no --- implies
that the information channel operates through portfolio choice rather than
through mechanical beta exposure: an MP shock moves discount rates, which
prices incorporate instantly, but gives investors no new reason to change
allocations; a CBI shock changes the expected distribution of returns across
classes, which is exactly what an allocation responds to. Third, the unwind
pattern suggests that part of the flow response reflects temporary
over-allocation --- consistent with slow-moving capital chasing news --- rather
than a permanent change in target allocations.

\section{Related literature}\label{sec:lit}

Our paper connects three literatures. The first is the measurement of
monetary-policy shocks. \cite{KUTTNER_2001} established the fed-futures
surprise methodology; \cite{GURKAYNAK_2005} decomposed announcements into
policy-path and timing components; \cite{JAROCINSKI_KARADI_2020} added the
sign-restriction separation of MP and CBI components, which we adopt in its
officially maintained form. \cite{BAUER_SWANSON_2023} showed that the
``Fed information effect'' survives --- and strengthens under --- controls for
the strong pre-FOMC evidence; our statement-based measures are constructed to
be comparable to theirs, and our flows specification includes their argument
as a design constraint.

The second is the real-economy and portfolio consequences of policy news.
\cite{KROENCKE_2024} document risk-on/risk-off shifts in fund flows around
FOMC meetings and are the closest paper to ours; our contribution relative to
theirs is the explicit MP-versus-CBI decomposition (they study the combined
surprise), the risk-ladder interaction, and the unwind result.
\cite{BANEGAS_2022} estimate flow responses to monetary-policy shocks by fund
strategy using a factor-augmented VAR shock; \cite{BROOKS_2018} document
bond-fund outflows after rate hikes with a lag. Neither separates the
information component, which is precisely where our results live.

The third is the measurement of statement tone. Dictionary scores
\citep{LOUGHRAN_MCDONALD_2011, BOUVARES_2020} and LLM-based scores
\citep{CURTI_KAZINNIK_2023, HANSEN_2025} both correlate with market reactions,
but they differ sharply in resolution across regimes. Our score-compression
result (\S\ref{sec:measurement}) --- the dictionary loses most of its
cross-meeting variance in the transition from forward guidance to
quantitative easing --- provides a measurement explanation for why tone effects
appear episodic in earlier work, and motivates the sentiment-versus-hawkishness
split that our price-side horse race exploits.

\section{Institutional background and identification}\label{sec:identification}

\subsection{Two kinds of news}\label{sec:twokinds}

Let $s_t$ denote the statement at announcement $t$, $i_t$ the high-frequency
announcement-window change in the policy rate, and $r_t$ the contemporaneous
stock-price change. The \jk{} restrictions impose that an MP shock moves $i_t$
and $r_t$ in opposite directions, while a CBI shock moves them together. The
estimator returns the two components $(mp_t, cbi_t)$ as linear combinations of
$(i_t, r_t)$ up to scale. We use the author-maintained update
(\texttt{jkshocks\_update\_fed}), taking the post-meeting-window series
(\texttt{MP\_pm}, \texttt{CBI\_pm}); every one of our 161 meeting dates
(2006--2026) matches a date in that file.

The interpretation convention we adopt throughout: a \emph{positive} CBI
shock is favourable economic news revealed by the statement (the professor's
``positive CBI = good economic news'' labelling), and a positive MP shock is a
contractionary policy surprise. Under this convention the flows results read
naturally: favourable information $\Rightarrow$ money moves toward risk.

\subsection{Why the statement matters for flows}\label{sec:whystatements}

Two features of the modern statement make it the relevant information vehicle
for fund managers. First, since 2005 the statement is a structured document
whose sentence-level content (assessment of activity, inflation, and the
balance of risks) maps into portfolio-relevant state variables. Second, the
manager's problem is an allocation problem: what matters is not the level of
any single price but the joint distribution across classes. A statement that
reveals ``the economy is stronger than we thought'' changes relative expected
returns --- bond classes lose, equity classes gain --- and therefore has a
first-order allocation implication, whereas a pure action surprise (already
priced in futures) changes conditional means only through the discount-rate
channel that no allocation strategy targets. Our hypothesis block formalises
this (\S\ref{sec:hypotheses}).

\section{Data}\label{sec:data}

\paragraph{Meetings and shocks.} 161 FOMC announcements, 2006--2026.
High-frequency surprises: fed-funds futures (target), Eurodollar/futures
(path), from the standard construction. Official MP/CBI: as above. Statement
text: body-only scrubbed text for all 161 statements; sentiment is
$0.5\times$LM $+\ 0.5\times$ central-bank word list per 100 words
(tone100); the LLM hawkishness score is the five-model consensus
($-100$ to $+100$) described in \S\ref{sec:measurement}.

\paragraph{Fund flows.} CRSP Mutual Fund Database, monthly, point-in-time
(fund header history matched at each fund-month). Classes are the seven
Morningstar-style groups above; class flows are TNA-weighted sums of
fund-level dollar flows divided by lagged class TNA
(\% of lagged TNA). Sample: 245 class-months, 1,099 panel observations with
controls (Tables use $n$ between 1,022 and 1,109 depending on the window and
the shock series' coverage).

\paragraph{Prices.} Same-day changes: CRSP value-weighted and equal-weighted
indices, S\&P 500, 10-year and 2-year Treasuries, 3-month T-bills.

\section{Hypotheses}\label{sec:hypotheses}

The theory block (rational-expectations inversion; asymmetric information
between the committee and the market; portfolio-choice under changing state
expectations) generates four testable implications. We state them in the
order the results section reports:

\begin{itemize}
\item \textbf{H1 (flows headline).} CBI shocks shift flows along the risk
ladder; MP shocks do not. (Formally: $\beta_{\text{Risk}\times\text{CBI}}>0$;
$\beta_{\text{Risk}\times\text{MP}}=0$.)
\item \textbf{H2 (safe asset).} Government-bond flows respond to CBI with the
opposite sign of the ladder result; the MP and CBI coefficients differ
(Wald).
\item \textbf{H3 (unwind).} The same-month flow response partially reverses
in the post-minus-same difference.
\item \textbf{H4 (breadth, not direction).} The FG-period price response loads
on sentiment breadth, not on the directional hawkishness score.
\end{itemize}

Note the deliberate numbering: flows are the paper's contribution and are
presented as H1--H3; the price-side tone results follow as H4. This inverts
the working-paper draft numbering and follows the structure recommended in
the November 2026 revision memo.

\section{Results I: fund flows}\label{sec:flows}

\subsection{Headline: the ladder result}\label{sec:ladder}

Table~\ref{tab:official} is the paper's core table. Columns 1--7 show the
H4 government-bond regressions (meeting-month cross-section, HC1); columns
8--11 show the H5 pooled risk-ladder panel (month-clustered). In the meeting
month, a favourable CBI shock pushes flows up the ladder ($+5.28$, $p=.005$)
and out of government bonds ($-43.1$, $p=.037$); the MP coefficient is
insignificant everywhere. One month later the CBI effects persist at roughly
$60$--$70\%$ of the initial magnitude (ladder $+3.45$, $p=.006$; gov-bond
$-24.5$, $p=.024$). By the difference window the ladder effect is statistically
exhausted ($-1.79$, $p=.37$).

The economic magnitude is sizeable. A one-standard-deviation favourable CBI
shock moves roughly $5.3$ percentage points of extra monthly flow toward a
one-step-riskier class relative to a one-step-safer one; at the average
government-bond class TNA of about \$110 billion this is on the order of
several billion dollars of reallocation per announcement.

\subsection{MP shocks: a disciplined null}\label{sec:mpnull}

The MP column of Table~\ref{tab:official} is a null we designed the paper to
have. Point estimates are small ($-0.41$ to $+0.42$ on the ladder) and
imprecisely signed across windows ($p=.47$--$.95$). This is not a power
failure: with the same panel, the CBI interaction is estimated at four
$t$-statistics. The null also makes economic sense --- an MP surprise moves
discount rates and prices immediately, but portfolio allocations target
expected relative returns, which an anticipated action does not change.

\subsection{The unwind}\label{sec:unwind}

The difference-window result (H3) quantifies the unwind: the same-month
ladder response partially reverses by month two (Table~\ref{tab:official},
diff row). Two interpretations are
consistent with the data: slow-moving capital overshoots on announcement and
corrects; or the same-month response reflects temporary liquidity provision
by flow-driven investors that mean-reverts. Distinguishing these is on the
agenda; the JIMF literature on delayed bond-fund responses
\citep{BROOKS_2018} suggests the first.

\section{Results II: tone and prices}\label{sec:prices}

\subsection{The FG-period tone effect}\label{sec:fgtone}

Over the full sample, statement sentiment adds little to returns beyond the
Target and Path surprises (Table~\ref{tab:horse}, full-sample rows:
$p=.25$--$.39$). In the forward-guidance regime --- 56 meetings, 2008--2015,
where the statement was the marginal instrument --- sentiment predicts
same-day Treasury and equity returns strongly: return on T, P, sentiment
gives $-1.01$ ($p=.010$, $R^2=32.8\%$) on CRSP VW, and the full-sample
interaction specification gives Tone$\times$FG $=-1.13$ ($p=.003$). An
orthogonalized version (sentiment residualised on Target and Path)
confirms: $-1.06$ ($p=.006$).

\subsection{Breadth, not direction (H4)}\label{sec:split}

Table~\ref{tab:split} runs the professor-requested split. In the FG
subsample, the LLM hawkishness score alone is insignificant
($p=.23$--$.33$); when both measures enter, sentiment dominates
($-1.11$, $p=.005$) and hawkishness remains marginal at best ($p=.081$).
The reading: prices respond to the \emph{richness} of the committee's
assessment across many dimensions, not to the single directional
hawkish/dovish axis --- consistent with the Bayesian-inversion view in which
the market updates on the whole state vector the statement reveals.

\subsection{The horse race with official shocks}\label{sec:horserace}

Table~\ref{tab:horse} (bottom rows) adds the official MP and CBI shocks to
the FG specification. Tone survives: $-0.96$ ($p=.009$) with official controls
versus $-1.22$ ($p=.001$) with sign-rule controls (the latter reproduces the
48\%-$R^2$ specification of the November deck). The official shock
coefficients themselves are marginally signed in this subsample
($p=.052$--$.078$), which is unsurprising given $n=56$ and the overlap between
announcement-window and same-day return measurement.

\section{Tone measurement: compression and the LLM}\label{sec:measurement}

Do our tone measures separate policy actions? Table~\ref{tab:measures} asks
which measures distinguish hikes from cuts. The Loughran--McDonald dictionary
does not ($p=.160$); the central-bank word list does ($p=.024$); their average
does ($p<.01$); the five-model-consensus LLM score separates them at
$p<10^{-9}$ with a $2.0$-SD gap. But the deeper finding is
\emph{compression}: across the 2022 regime transition the dictionary's
cross-meeting standard deviation collapses while the LLM score's does not ---
the dictionary loses roughly two orders of magnitude of discriminating
variance exactly where the policy regime changes. This measurement failure
explains why dictionary-based tone effects look episodic in the literature,
and is the practical argument for LLM-based scoring in this literature.

\section{Robustness}\label{sec:robustness}

\paragraph{Inference.} Table~\ref{tab:se} reports the six-covariance-estimator
grid for the headline sign-rule regressions. The FG-period tone result
survives everything (classical/HC1/NW/wild: $p=.003$--$.012$); the sign-rule
flow results do not (HC1 $p=.34$--$.43$, cluster $p=.47$). This is why the
paper's primary flow specifications use the official series.

\paragraph{Placebo and power.} The FG interaction is not an artifact of
arbitrary period choice: across 95 consecutive 56-meeting windows, the actual
estimate is the most negative ($p=.01$ share-level). The minimum detectable
effect at 80\% power for the FG tone result ($1.07$) brackets the estimate
($1.01$): the design is powered for the effect it finds.

\paragraph{Multiple testing.} Of 42 sign-rule H4 coefficients, none survives
Holm or Benjamini--Hochberg correction under either classical or HC1
standard errors (Table~\ref{tab:multiple}). We flag this honestly: the
sign-rule flow grid is dead as an evidentiary basis, and the paper's flow
claims rest exclusively on the official-series specifications of
Table~\ref{tab:official} --- which are three pre-registered, theory-derived
coefficients, not a grid.

\paragraph{Beta measurement.} Replacing the risk rank with downside/upside
beta splits, continuous betas, and rolling betas leaves every interaction
insignificant under the sign-rule series (the measurement is not the
problem --- the series is). State dependence on credit-spread and inflation
conditions: none ($p>.22$ everywhere).

\paragraph{Data integrity.} A classification bug was found and fixed during
the audit: meetings with missing same-day returns had been coded as MP shocks
by the sign rule ($91\to83$ MP meetings). The correction changes no flow
result materially (gov-bond CBI $+0.76$, $p=.043$ before and after). All
results on the official series are unaffected by construction.

\section{Conclusion}\label{sec:conclusion}

Federal Reserve statements do two things: they act, and they reveal. The
prices literature has established that markets react to both. This paper
establishes that \emph{portfolios} respond to the revelation, not the action:
information shocks move money up and down the risk ladder with
month-clustered significance, conventional shocks do not move it at all, and
the response partially unwinds --- portfolios, like prices, move first on the
day and digest later. The statement's sentiment breadth, not its directional
tone, is what prices price. And the identification exercise delivers a
methodological warning: statement-text proxies for information shocks can be
nearly uncorrelated with the high-frequency object they seek to capture, and
inference that survives under one can die under the other.

Three extensions follow directly. Retail-versus-institutional splits will
identify who moves; the SPF forecast-error version of CBI will provide a
non-financial validation of the information component; and the unwind result
invites a quantitative model of slow-moving capital under information shocks.

\bibliographystyle{plainnat}
\begin{thebibliography}{20}

\bibitem[Banegas et al.(2022)]{BANEGAS_2022}
Banegas, A., Montes-Rojas, G., and Siga, L. (2022). The effects of U.S.
monetary policy shocks on mutual fund investing. \emph{Journal of International
Money and Finance}.

\bibitem[Bauer and Swanson(2023)]{BAUER_SWANSON_2023}
Bauer, M. D., and Swanson, E. T. (2023). An alternative explanation for the
``Fed information effect.'' \emph{Journal of Monetary Economics}.

\bibitem[Bernanke and Kuttner(2005)]{BERNANKE_KUTTNER_2005}
Bernanke, B. S., and Kuttner, K. N. (2005). What explains the stock market's
reaction to Federal Reserve policy? \emph{Journal of Finance}, 60(3),
1221--1257.

\bibitem[Bouvard et al.(2020)]{BOUVARES_2020}
Bouvard, M., Chaigneau, P., and De Motta, A. (2020). Central bank word lists.
(Unpublished manuscript.)

\bibitem[Brooks et al.(2018)]{BROOKS_2018}
Brooks, J., Katz, M., and Lustig, H. (2018). Post-FOMC announcement drift in
U.S. bond markets. NBER Working Paper 25127.

\bibitem[Curti and Kazinnik(2023)]{CURTI_KAZINNIK_2023}
Curti, C., and Kazinnik, S. (2023). Let's face it: Quantifying the impact of
FOMC communication. \emph{Journal of Financial Economics}.

\bibitem[G\"urkaynak et al.(2005)]{GURKAYNAK_2005}
G\"urkaynak, R. S., Sack, B., and Swanson, E. T. (2005). Do actions speak
louder than words? \emph{International Journal of Central Banking}.

\bibitem[Hansen et al.(2025)]{HANSEN_2025}
Hansen, A., et al. (2025). Large language models for central-bank
communication analysis. (Unpublished manuscript --- verify authors before
submission.)

\bibitem[Jaroci\'nski and Karadi(2020)]{JAROCINSKI_KARADI_2020}
Jaroci\'{n}ski, M., and Karadi, P. (2020). Deconstructing monetary policy
surprises. \emph{American Economic Journal: Macroeconomics}, 12(2), 26--59.

\bibitem[Kroencke et al.(2024)]{KROENCKE_2024}
Kroencke, T. A., Schmeling, M., and Schrimpf, A. (2024). The FOMC risk shift.
CEPR Discussion Paper 14037.

\bibitem[Kuttner(2001)]{KUTTNER_2001}
Kuttner, K. N. (2001). Monetary policy surprises and interest rates: Evidence
from the Fed funds futures market. \emph{Journal of Monetary Economics},
47(3), 523--544.

\bibitem[Loughran and McDonald(2011)]{LOUGHRAN_MCDONALD_2011}
Loughran, T., and McDonald, B. (2011). When is a liability not a liability?
\emph{Journal of Finance}, 66(1), 35--65.

\end{thebibliography}

\newpage
\appendix

\section*{Tables}\addcontentsline{toc}{section}{Tables}

\begin{table}[H]
\centering
\begin{threeparttable}
\caption{Official \jk{} series: fund flows by window (the paper's core table)}
\label{tab:official}
% --- auto-generated by gen_tables.py; do not edit numbers by hand ---
\begin{tabular}{lrccccccccc}
\toprule
 & & \multicolumn{5}{c}{H4: Gov-bond flows} & \multicolumn{4}{c}{H5: Risk ladder} \\
\cmidrule{3-7}\cmidrule{8-11}
Window & N & CBI $\beta$ & (SE) & MP $\beta$ & (SE) & Wald & Risk$\times$CBI & (SE) & Risk$\times$MP & (SE) \\
\midrule
Same month & 157 & -43.061$^{**}$ & (20.438) & +5.005 & (6.899) & .025 & +5.277$^{***}$ & (1.858) & -0.414 & (0.711) \\
Post month & 158 & -24.537$^{**}$ & (10.738) & +0.740 & (5.082) & .031 & +3.449$^{***}$ & (1.238) & -0.030 & (0.517) \\
Diff (post $-$ same) & 157 & +17.305 & (19.189) & -4.614 & (4.786) & .272 & -1.787 & (1.980) & +0.417 & (0.485) \\
\bottomrule
\end{tabular}
\begin{tablenotes}\footnotesize
\item \textbf{Note.} Shocks are the official Jaroci\"{\i}n--Karadi (2020) decomposition (updated series, MP\_pm/CBI\_pm), matched to 161/161 meeting dates. H4: government-bond fund flows on MP and CBI, meeting-month cross-section, HC1 standard errors; Wald tests $\beta_{MP}=\beta_{CBI}$. H5: class-month panel with Risk$\times$MP and Risk$\times$CBI interactions and controls (log TNA, 12-month flow volatility), standard errors clustered by meeting month. Diff window = post-month minus same-month flow. $^{*}$p$<$.10, $^{**}$p$<$.05, $^{***}$p$<$.01. Source: robustness\_tests/results.json, test2.
\end{tablenotes}
\end{threeparttable}
\end{table}

\begin{table}[H]
\centering
\begin{threeparttable}
\caption{Shock series comparison and the classification fix}
\label{tab:series}
% --- auto-generated; do not edit numbers by hand ---
\begin{tabular}{lcccc}
\toprule
 & MP meetings & CBI meetings & Gov CBI $\beta$ (same) & Risk$\times$CBI $\beta$ (same) \\
\midrule
As published (2025--26 coded MP) & 91 & 60 & +0.759$^{**}$ & -0.066$^{**}$ \\
Classification fix (valid-return gate) & 83 & 60 & +0.763$^{**}$ & -0.066$^{**}$ \\
Official JK series (clustered) & 161 & 161 & -43.061$^{**}$ & +5.277$^{***}$ \\
\bottomrule
\end{tabular}
\begin{tablenotes}\footnotesize
\item \textbf{Note.} Row 1: sign-rule classification as published (2025--26 meetings without same-day returns coded as MP). Row 2: meetings without valid same-day CRSP returns excluded (fix of the NaN-sign classification bug; 91 $\to$ 83 MP meetings). Row 3: official JK shocks, month-clustered SEs. Correlations between series: MP $\rho = 0.529$; CBI $\rho = 0.038$; within-official MP--CBI $\rho = -0.009$. Source: baseline\_check + test2.
\end{tablenotes}
\end{threeparttable}
\end{table}

\begin{table}[H]
\centering
\begin{threeparttable}
\caption{Tone and prices: FG-period horse race with official MP/CBI controls}
\label{tab:horse}
% --- auto-generated ---
\begin{tabular}{lrccccccccccc}
\toprule
Specification & N & R$^2$ & \multicolumn{2}{c}{Tone} & \multicolumn{2}{c}{MP} & \multicolumn{2}{c}{CBI} & \multicolumn{2}{c}{Target} & \multicolumn{2}{c}{Path} \\
\cmidrule(l){4-5}\cmidrule(l){6-7}\cmidrule(l){8-9}\cmidrule(l){10-11}\cmidrule(l){12-13}
 & & & $\beta$ & (SE) & $\beta$ & (SE) & $\beta$ & (SE) & $\beta$ & (SE) & $\beta$ & (SE) \\
\midrule
T + P + tone + MP(sign rule) + CBI(sign rule) & 56 & 50.5\% & -1.218$^{***}$ & (0.334) & -1.344$^{***}$ & (0.325) & +1.633$^{***}$ & (0.598) & +0.289 & (0.310) & -0.346 & (0.240) \\
T + P + tone + MP(Official JK) + CBI(Official JK) & 56 & 38.8\% & -0.959$^{***}$ & (0.353) & -89.208$^{**}$ & (38.681) & -107.081$^{*}$ & (55.572) & +2.560$^{*}$ & (1.425) & +3.013$^{*}$ & (1.514) \\
tone + MP(Official JK) + CBI(Official JK) & 56 & 35.0\% & -0.999$^{***}$ & (0.358) & -14.648$^{***}$ & (4.073) & -32.161 & (38.222) & --- & --- \\
tone + MP(Official JK) + CBI(Official JK) & 150 & 8.0\% & -0.135 & (0.126) & -7.427$^{**}$ & (3.455) & -3.277 & (9.811) & --- & --- \\
\bottomrule
\end{tabular}
\begin{tablenotes}\footnotesize
\item \textbf{Note.} Dependent variable: same-day CRSP value-weighted return, percentage points. Rows 1--3: forward-guidance subsample (56 meetings, dataset regime flag). Row 4: full sample (150 meetings with returns). HC1 standard errors. Row 1 reproduces the published horse-race specification; rows 2--3 replace the sign-rule shocks with the official JK series; official-series coefficients are in the series' native units, so their scale is not comparable across rows --- read the $p$-values. Source: professor\_revisions.json.
\end{tablenotes}
\end{threeparttable}
\end{table}

\begin{table}[H]
\centering
\begin{threeparttable}
\caption{Sentiment vs.\ hawkishness (professor memo E11)}
\label{tab:split}
% --- auto-generated ---
\begin{tabular}{lrccccc}
\toprule
 & & & \multicolumn{2}{c}{Sentiment (tone100)} & \multicolumn{2}{c}{Hawkishness (LLM)} \\
\cmidrule(l){4-5}\cmidrule(l){6-7}
Specification & N & R$^2$ & $\beta$ & SE & $\beta$ & SE \\
\midrule
FG subsample: T + P + sentiment & 56 & 32.8\% & -1.013$^{**}$ & (0.381) &  &  \\
FG subsample: T + P + hawkishness & 56 & 23.6\% &  &  & -0.006 & (0.005) \\
FG subsample: sentiment + hawkishness & 56 & 36.9\% & -1.113$^{***}$ & (0.379) & -0.008$^{*}$ & (0.004) \\
Full sample --- return on T, P, sentiment & 150 & 6.3\% & -0.145 & (0.125) &  &  \\
full: R ~ T + P + hawkishness & 150 & 6.5\% &  &  & -0.003 & (0.003) \\
full: sentiment + hawkishness together & 150 & 6.8\% & -0.103 & (0.119) & -0.002 & (0.003) \\
\bottomrule
\end{tabular}
\begin{tablenotes}\footnotesize
\item \textbf{Note.} Sentiment = 0.5$\cdot$LM + 0.5$\cdot$central-bank word list (tone100). Hawkishness = five-model-consensus LLM score ($-100$ to $+100$). HC1 standard errors. FG subsample uses the dataset's regime flag (56 meetings). Source: professor\_revisions.json.
\end{tablenotes}
\end{threeparttable}
\end{table}

\begin{table}[H]
\centering
\begin{threeparttable}
\caption{Sign-rule robustness: six covariance estimators}
\label{tab:se}
% --- auto-generated ---
\begin{tabular}{lrcccccc}
\toprule
Result & $\beta$ & Classical & HC1 & Newey--West & Month cluster & Two-way & Wild bootstrap \\
\midrule
Return on T, P, tone (full) & -0.145 & .260 & .248 & .277 & --- & --- & .248 \\
\quad + FG interaction: Tone$\times$FG & -1.133 & .003 & .003 & .005 & --- & --- & .003 \\
FG subsample: return on T, P, tone & -1.013 & .005 & .010 & .011 & --- & --- & .012 \\
Gov-bond flow on CBI (same month) & +0.763 & .043 & .342 & .341 & --- & --- & --- \\
Risk$\times$CBI, class-month panel (same) & -0.066 & .041 & .434 & --- & .466 & --- & --- \\
\bottomrule
\end{tabular}
\begin{tablenotes}\footnotesize
\item \textbf{Note.} $p$-values under six covariance estimators for the five headline sign-rule regressions; column 2 of each row is the classical-SE point estimate of the tone (rows 1--3) or shock (rows 4--5) coefficient. --- = estimator not applicable (cross-section has no cluster dimension). The flow results (last two rows) do not survive robust inference; the FG-period tone result does. Source: robustness\_tests/results.json, test1.
\end{tablenotes}
\end{threeparttable}
\end{table}

\begin{table}[H]
\centering
\begin{threeparttable}
\caption{Do tone measures separate hikes from cuts?}
\label{tab:measures}
% --- auto-generated ---
\begin{tabular}{lccccc}
\toprule
Measure & Mean (hikes) & Mean (cuts) & Mean (holds) & Hike--cut gap (SD) & $p$ \\
\midrule
Loughran--McDonald dictionary & 3.59 & 3.03 & 4.58 & 0.39 & .160 \\
Central-bank word list & 2.69 & 1.62 & 0.98 & 0.67$^{**}$ & .024 \\
Combined (0.5 LM + 0.5 CB) & 3.14 & 2.33 & 2.78 & 1.01$^{***}$ & .002 \\
LLM hawkishness ($-100$ to $+100$) & 51.67 & -43.12 & -15.93 & 2.00$^{***}$ & $<$ .001 \\
\bottomrule
\end{tabular}
\begin{tablenotes}\footnotesize
\item \textbf{Note.} Mean score by decision (21 hikes, 16 cuts, 124 holds). Gap in within-measurement standard-deviation units; $p$-value from a two-sample test of hike vs.\ cut means. Source: tables\_all\_results.xlsx, Table 8 (tone\_measures\_current.json).
\end{tablenotes}
\end{threeparttable}
\end{table}

\begin{table}[H]
\centering
\begin{threeparttable}
\caption{Power, placebo, and multiple testing}
\label{tab:multiple}
% --- auto-generated ---
\begin{tabular}{lcccc}
\toprule
 & Tests & Raw $p<.05$ & Survive Holm & Survive BH \\
\midrule
Classical SE & 42 & 6 & 0 & 0 \\
HC1 SE & 42 & 2 & 0 & 0 \\
\bottomrule
\end{tabular}

\vspace{0.6em}

\begin{tabular}{lccc}
\toprule
Power / placebo & Estimate & SE & MDE / $p$ \\
\midrule
Tone on returns, FG only & -1.013 & 0.381 & MDE 1.07 \\
Placebo: random 56-meeting windows & -1.133 & & $p = 0.0$; 95 blocks, 1.1\% at least as negative \\
\bottomrule
\end{tabular}
\begin{tablenotes}\footnotesize
\item \textbf{Note.} Top: 42 H4 sign-rule coefficient tests (7 classes $\times$ 3 windows $\times$ 2 shocks), Holm and Benjamini--Hochberg corrections. Bottom: minimum detectable effect at 80\% power for the FG-period tone result, and the placebo distribution over consecutive 56-meeting windows. Source: robustness\_tests/results.json, test6.
\end{tablenotes}
\end{threeparttable}
\end{table}

\begin{thebibliography}{40}
\setlength{\itemsep}{0pt}

\bibitem[Banegas et al.(2022)]{BANEGAS_2022}
Banegas, A., Montes-Rojas, G., and Siga, L. (2022). ``The Effects of U.S.
Monetary Policy Shocks on Mutual Fund Investing.'' \emph{Journal of
International Money and Finance}.

\bibitem[Bauer and Swanson(2023)]{BAUER_SWANSON_2023}
Bauer, M. D., and Swanson, E. T. (2023). ``An Alternative Explanation for the
`Fed Information Effect.''' \emph{Journal of Monetary Economics}.

\bibitem[Bernanke and Kuttner(2005)]{BERNANKE_KUTTNER_2005}
Bernanke, B. S., and Kuttner, K. N. (2005). ``What Explains the Stock Market's
Reaction to Federal Reserve Policy?'' \emph{Journal of Finance} 60(3),
1221--1257.

\bibitem[Bouvard et al.(2020)]{BOUVARES_2020}
Bouvard, M., Chaigneau, P., and De Motta, A. (2020). ``Transparency in the
Financial System.'' \emph{Journal of Finance} (statement-transparency
mechanisms).

\bibitem[Brooks et al.(2018)]{BROOKS_2018}
Brooks, J., Katz, M., and Lustig, H. (2018). ``Post-FOMC Announcement Drift in
U.S. Bond Markets.'' NBER Working Paper 25127.

\bibitem[Curti and Kazinnik(2023)]{CURTI_KAZINNIK_2023}
Curti, C., and Kazinnik, S. (2023). ``Let's Face It: Quantifying Investor
Sentiment from the Visual Dimension of the Federal Reserve Press
Conferences.'' \emph{Journal of Financial Economics}.

\bibitem[G\"urkaynak et al.(2005)]{GURKAYNAK_2005}
G\"urkaynak, R. S., Sack, B., and Swanson, E. T. (2005). ``Do Actions Speak
Louder Than Words? The Response of Asset Prices to Monetary Policy Actions and
Statements.'' \emph{International Journal of Central Banking} 1(1), 55--93.

\bibitem[Hansen(2025)]{HANSEN_2025}
Hansen, S., McMahon, M., et al. (2025). Large-language-model measurement of
central-bank communication. (Concurrent LLM-scoring work.)

\bibitem[Jaroci\'nski(2024)]{JAROCINSKI_2024}
Jaroci\'nski, M. (2024). ``Estimating the Fed's Unconventional Policy Shocks.''
\emph{Journal of Monetary Economics}. (Replication files:
\texttt{github.com/marekjarocinski/jkshocks\_update\_fed}.)

\bibitem[Jaroci\'nski and Karadi(2020)]{JAROCINSKI_KARADI_2020}
Jaroci\'nski, M., and Karadi, P. (2020). ``Deconstructing Monetary Policy
Surprises --- The Roles of Information Shocks.'' \emph{Journal of Monetary
Economics} 113, 48--67.

\bibitem[Kroencke et al.(2024)]{KROENCKE_2024}
Kroencke, T. A., Schmeling, M., and Schrimpf, A. (2024). ``The FOMC Risk
Shift.'' CEPR Discussion Paper DP14037.

\bibitem[Kuttner(2001)]{KUTTNER_2001}
Kuttner, K. N. (2001). ``Monetary Policy Surprises and Interest Rates: Evidence
from the Fed Funds Futures Market.'' \emph{Journal of Monetary Economics}
47(3), 523--544.

\bibitem[Loughran and McDonald(2011)]{LOUGHRAN_MCDONALD_2011}
Loughran, T., and McDonald, B. (2011). ``When Is a Liability Not a Liability?
Textual Analysis, Dictionaries, and 10-Ks.'' \emph{Journal of Finance} 66(1),
35--65.

\end{thebibliography}

\end{document}
